% HRMA Kullanma Kılavuzu (Türkçe) — xelatex ile derlenir
% Derleme: xelatex HRMA-Kullanma-Kilavuzu-TR.tex (iki kez, içindekiler için)
\documentclass[11pt,a4paper]{article}

\usepackage{fontspec}
\usepackage[turkish]{babel}
\usepackage{geometry}
\geometry{a4paper, margin=2.4cm, top=2.6cm, bottom=2.6cm}
\usepackage{graphicx}
\usepackage{xcolor}
\usepackage{fancyhdr}
\usepackage{titlesec}
\usepackage{enumitem}
\usepackage{tcolorbox}
\usepackage{booktabs}
\usepackage{array}
\usepackage{hyperref}
\usepackage{parskip}
\tcbuselibrary{skins,breakable}

% ---- Renk paleti (uygulamanın koyu temasıyla uyumlu) ----
\definecolor{hrmacyan}{HTML}{0891A8}
\definecolor{hrmadark}{HTML}{0A1524}
\definecolor{hrmaink}{HTML}{1B2A3A}
\definecolor{hrmagray}{HTML}{5A6B7A}
\definecolor{hrmapanel}{HTML}{EEF4F8}
\definecolor{hrmawarn}{HTML}{B26A00}

\hypersetup{
  colorlinks=true, linkcolor=hrmacyan, urlcolor=hrmacyan,
  pdftitle={HRMA Kullanma Kılavuzu}, pdfauthor={Berke Tezgöçen}
}

% ---- Başlık biçimleri ----
\titleformat{\section}{\Large\bfseries\color{hrmacyan}}{\thesection.}{0.5em}{}
\titleformat{\subsection}{\large\bfseries\color{hrmadark}}{\thesubsection}{0.5em}{}
\titlespacing*{\section}{0pt}{1.4em}{0.6em}

% ---- Üst/alt bilgi ----
\pagestyle{fancy}
\fancyhf{}
\fancyhead[L]{\small\color{hrmagray}HRMA — Roket Motoru Analiz Paketi}
\fancyhead[R]{\small\color{hrmagray}Kullanma Kılavuzu}
\fancyfoot[C]{\small\color{hrmagray}\thepage}
\renewcommand{\headrulewidth}{0.4pt}
\renewcommand{\footrulewidth}{0pt}

% ---- Kutu stilleri ----
\newtcolorbox{notkutu}{
  colback=hrmapanel, colframe=hrmacyan, boxrule=0.6pt, arc=3pt,
  left=8pt, right=8pt, top=6pt, bottom=6pt, breakable}
\newtcolorbox{uyarikutu}{
  colback=yellow!8, colframe=hrmawarn, boxrule=0.8pt, arc=3pt,
  left=8pt, right=8pt, top=6pt, bottom=6pt, breakable}

% Görsel + açıklama yardımcı komutu
\newcommand{\gorsel}[2]{%
  \begin{center}
  \fbox{\includegraphics[width=0.96\linewidth]{#1}}\\[3pt]
  {\small\color{hrmagray}\itshape #2}
  \end{center}}

\begin{document}
% Türkçe babel '=' ':' '!' karakterlerini kısayol yapıyor; bu, enumitem/tcolorbox
% gibi paketlerin anahtar=değer sözdizimini bozuyor. Kısayolları kapatıyoruz.
\shorthandoff{=:!}

% ================= KAPAK =================
\begin{titlepage}
\centering
\vspace*{2.5cm}
{\fontsize{54}{60}\selectfont\bfseries\color{hrmadark}HRMA}\\[4pt]
{\Large\color{hrmacyan}Yüksek Doğruluklu Roket Motoru Analiz ve Tasarım Paketi}\\[6pt]
{\large\color{hrmagray}Hibrit • Katı • Sıvı Roket Motorları}\\[2.2cm]

{\Huge\bfseries\color{hrmadark}Kullanma Kılavuzu}\\[10pt]
{\large\color{hrmagray}Örnekli, adım adım anlatım}\\[3cm]

{\large Sürüm 2.6.1 \quad|\quad Temmuz 2026}\\[6pt]
{\small\color{hrmagray}Geliştiren: Berke Tezgöçen \quad•\quad Fikir \& Test: Ayberk Cem Aksoy}\\[1.5cm]

\begin{notkutu}
\small\color{hrmadark}
\textbf{Kapsam.} HRMA; ön tasarım ve eğitim amaçlı, kapalı-biçim ve
bir-boyutlu mühendislik bağıntılarına dayanan bir çözümleme aracıdır.
Uçuşa elverişlilik (flight qualification) aracı \emph{değildir}. Sonuçlar
uçuş öncesi bağımsız bir kodla (CEA / RPA / openMotor) ve fiziksel testle
doğrulanmalıdır. Ayrıntı için kılavuzun sonundaki ``Kapsam ve Sınırlar''
bölümüne bakın.
\end{notkutu}
\end{titlepage}

% ================= İÇİNDEKİLER =================
\tableofcontents
\newpage

% ================= 1. GİRİŞ =================
\section{HRMA Nedir?}

HRMA (Roket Motoru Analiz Paketi), hibrit, katı ve sıvı roket motorlarının
ön tasarımı ve çözümlemesi için geliştirilmiş bir masaüstü uygulamasıdır.
Uygulama, yerel bir hesaplama motorunu (Flask) kendi penceresinde çalıştırır
ve sonuçları koyu temalı, etkileşimli bir arayüzde sunar: gerçek çözücü
çıktısından üretilen 3B motor modeli, mühendislik panellerinden oluşan bir
Analiz Güvertesi ve çalışan CAD / teknik resim / rapor dışa aktarımları.

\textbf{Kimler için?} Roket takımları, lisans/lisansüstü öğrenciler,
öğretim üyeleri ve motor fiziğini uçtan uca görmek isteyen herkes. Program
size bir motoru \emph{sıfırdan} tasarlatmaz; verdiğiniz parametrelerden
geometri, performans ve uyarılar üretir; siz de bunları mühendislik
muhakemesiyle değerlendirirsiniz.

\begin{notkutu}
\small\color{hrmadark}
\textbf{Her şey yerelde ve çevrimdışı çalışır.} Girdiğiniz hiçbir tasarım
verisi bilgisayarınızdan çıkmaz. Programın internete eriştiği tek an,
açılışta yeni sürüm olup olmadığını GitHub'dan sorduğu andır; bu da
başarısız olursa program normal çalışmaya devam eder.
\end{notkutu}

\subsection{Bu kılavuz nasıl kullanılır?}

Kılavuz, tek bir \textbf{örnek motor} üzerinden ilerler: HTPB yakıtı ve sıvı
N\textsubscript{2}O oksitleyici kullanan, 1000~N itkili, 10~s yanma süreli
bir hibrit motor. Aynı adımlar katı ve sıvı motorlar için de geçerlidir;
farklar ilgili bölümlerde belirtilmiştir. Ekran görüntüleri doğrudan
uygulamadan alınmıştır.

% ================= 2. KURULUM =================
\section{Kurulum ve İlk Açılış}

\subsection{Kurulum}

\begin{itemize}[leftmargin=1.4em]
  \item \textbf{Windows 10/11:} \texttt{HRMA-Setup-2.6.1.exe} dosyasını
  indirin, çift tıklayın ve sihirbazı izleyin (İleri, İleri, Kur). Kurulum
  kullanıcı bazlıdır; yönetici hakkı gerekmez. Windows SmartScreen imzasız
  olduğu için uyarabilir: ``Ek bilgi'' $\rightarrow$ ``Yine de çalıştır''.
  \item \textbf{macOS 11+ (Apple Silicon):} \texttt{HRMA-Setup-2.6.1-macOS.dmg}
  dosyasını açın ve \texttt{HRMA} simgesini \texttt{Applications} klasörüne
  sürükleyin. İlk açılışta uygulamaya sağ tıklayıp ``Aç'' deyin (imzasız
  uygulamalar için Gatekeeper bunu bir kez ister).
\end{itemize}

Python ve tüm kütüphaneler paketin içindedir; ayrıca bir kurulum gerekmez.

\subsection{İlk açılış}

Uygulama, hesaplama motorları arka planda yüklenirken yaklaşık bir saniye
içinde bir açılış ekranı gösterir, ardından ana sayfayı kendi penceresinde
açar. Ana sayfada üç motor tipi kartı ve formüller sayfasına bir bağlantı
bulunur.

\gorsel{img/tr-00-home.png}{Ana sayfa: üç motor tipini (Hibrit, Katı, Sıvı)
seçebileceğiniz kartlar ve üst şeritte dil seçici, Sürüm Notları, Kullanma
Kılavuzu ve Ayarlar bağlantıları.}

\begin{notkutu}
\small\color{hrmadark}
\textbf{Dil.} Sağ üstteki seçiciden arayüzü Türkçe veya İngilizce
yapabilirsiniz; seçim tüm sayfalarda ve grafik metinlerinde geçerlidir.
Bu kılavuzun İngilizce sürümü de uygulamayla birlikte gelir.
\end{notkutu}

% ================= 3. ARAYÜZ =================
\section{Ana Sayfalar}

Uygulamanın beş ana sayfası vardır:

\begin{center}
\begin{tabular}{@{}>{\raggedright}p{3.2cm} p{2.2cm} p{8.5cm}@{}}
\toprule
\textbf{Sayfa} & \textbf{Yol} & \textbf{Amaç} \\
\midrule
Ana Sayfa & \texttt{/} & Motor tipi seçimi ve formül başvurusu \\
Hibrit Tasarımcı & \texttt{/hybrid} & Tam hibrit motor tasarımı (en kapsamlı sayfa) \\
Katı Tasarımcı & \texttt{/solid} & Grain geometrisi, balistik, Monte Carlo \\
Sıvı Tasarımcı & \texttt{/liquid} & Besleme sistemi, enjektör, rejeneratif soğutma, çevrim \\
Formüller & \texttt{/formulas} & Çözücülerin kullandığı denklemler (MathJax) \\
\bottomrule
\end{tabular}
\end{center}

Üç tasarımcı sayfası da aynı yapıyı paylaşır: bir girdi formu, bir Hesapla
düğmesi, sonuç göstergeleri (HUD), bağımsız paneller (geçici rejim /
enjektör / 6-SD), Analiz Güvertesi, 3B model ve dışa aktarma düğmeleri.

% ================= 4. ÖRNEK =================
\section{Örnek: Bir Hibrit Motor Tasarlayalım}

Bu bölümde baştan sona bir örnek yapacağız. \texttt{/hybrid} sayfasını açın.

\subsection{1. Adım — Parametreleri girin}

Form yukarıdan aşağıya doğru mantıklı varsayılanlarla gelir. Örneğimiz için
şu değerleri kullanıyoruz (çoğu zaten varsayılan):

\begin{center}
\begin{tabular}{@{}p{5.5cm} p{4cm} p{4.5cm}@{}}
\toprule
\textbf{Alan} & \textbf{Değer} & \textbf{Açıklama} \\
\midrule
Motor Adı & \texttt{Ornek-H1} & Rapor ve dışa aktarmalarda kullanılır \\
İrtifa & \texttt{0 m} & Deniz seviyesi (ortam basıncı otomatik) \\
İtki & \texttt{1000 N} & Hedeflenen itki düzeyi \\
Yanma Süresi & \texttt{10 s} & Toplam impuls buradan türetilir \\
O/F Oranı & \texttt{2.5} & Oksitleyici/yakıt kütle oranı \\
Oda Basıncı & \texttt{20 bar} & Yanma odası basıncı \\
Tank Basıncı & \texttt{50 bar} & Oda basıncını yeterince aşmalı \\
Yakıt & \texttt{HTPB} & Regresyon katsayıları otomatik gelir \\
Oksitleyici & \texttt{N\textsubscript{2}O (sıvı)} & Yoğunluk sıcaklıkla güncellenir \\
Enjektör & \texttt{Showerhead} & Tip başına parametreler \\
\bottomrule
\end{tabular}
\end{center}

\gorsel{img/tr-01-form.png}{Girdi formu. ``Optimumu Bul'' düğmesi O/F oranını
c* açısından en iyi değere tarar. Her alandaki ``?'' simgesi kısa bir
açıklama gösterir.}

\begin{notkutu}
\small\color{hrmadark}
\textbf{0 = otomatik.} Genişleme oranı, oda çapı, port çapı gibi alanlara
\texttt{0} girerseniz çözücü o boyutu kendisi belirler. Kendi verinizi
(örneğin statik ateşten türettiğiniz regresyon katsayıları) girerek
varsayılanları ezebilirsiniz.
\end{notkutu}

\subsection{2. Adım — Hesaplayın}

\textbf{Hesapla} düğmesine basın. Çözücü, formu \texttt{/calculate} ucuna
gönderir ve birkaç saniye içinde sonuçları doldurur.

\gorsel{img/tr-02-results.png}{Sonuç göstergeleri (Performans Özeti),
Uyarılar paneli ve tasarım tabloları. Örneğimizde Özgül İtki 185.6~s,
İtki 1000~N, Oda Basıncı 20~bar, Kütle Debisi 0.549~kg/s ve Yapısal
Asgari Emniyet Katsayısı 2.89 çıktı.}

\textbf{Sonuçları okurken:}

\begin{itemize}[leftmargin=1.4em]
  \item \textbf{Performans Özeti (HUD):} Manşet sayılar — itki, toplam
  impuls, Isp, c*, CF, oda basıncı, boğaz ve çıkış çapları, kütle debileri
  ve bir emniyet durumu rozeti.
  \item \textbf{Uyarılar paneli:} Çözücüden ve doğrulama sisteminden gelen
  tasarım-ölçütü mesajları. \textbf{Sayılara güvenmeden önce bunları
  okuyun.} Örneğimizde üç uyarı var: c* normal aralığın hemen altında
  ($1325$ m/s), genişleme oranı düşük ($\varepsilon = 3.4$) ve ``flash
  boiling'' (ani kaynama) riski. Bunlar hata değil; tasarımın hangi
  varsayımlarla sınırda olduğunu söyler.
  \item \textbf{Tasarım tabloları:} Tam geometri — lüle kontur ölçüleri ve
  açıları, grain/port geometrisi, enjektör planı, cidar kalınlığı.
\end{itemize}

\begin{uyarikutu}
\small\color{hrmadark}
\textbf{Uyarılar bir özelliktir, kusur değil.} HRMA, hesaplayamadığı bir
şeye sayı uydurmak yerine açıkça uyarır veya ``modellenmedi'' der.
Kırmızı/turuncu bir uyarı gördüğünüzde, o parametreyi el hesabıyla veya
bağımsız bir kodla çapraz kontrol edin.
\end{uyarikutu}

\subsection{3. Adım — 3B modeli inceleyin}

Sonuçların altında, çözücünün ürettiği geometriden canlı kurulan bir 3B
motor modeli bulunur. Kesit, ölçüler, patlatılmış görünüm, ısı haritası ve
yanma animasyonu düğmeleriyle modeli inceleyebilirsiniz.

\gorsel{img/tr-03-3d.png}{3B motor simülasyonu (kesit görünümü). Ölçüler
doğrudan çözücü çıktısından gelir: oda çapı 83.5~mm, boğaz çapı 21.7~mm,
çıkış çapı 40.2~mm, toplam uzunluk 1377~mm. Alt şeritteki zaman çubuğu
yanma boyunca port çapının nasıl büyüdüğünü canlandırır.}

\begin{notkutu}
\small\color{hrmadark}
\textbf{Etkileşimli tasarım modu.} Bir hesaplamadan sonra oda çapı, L* ve
genişleme oranı kaydırıcıları geometriyi yaklaşık bir saniyede yeniden
hesaplar; 3B model ve kesit görünümü, tam bir yeniden hesaplama olmadan
anında güncellenir.
\end{notkutu}

\subsection{4. Adım — Grafikler ve kesit}

Performans grafikleri (itki eğrisi, basınç, regresyon hızı ve port büyümesi)
uygulamanın koyu temasında ve tamamen çevrimdışı çizilir. Kesit görünümü,
çözücünün kullandığı geometriden üretilen ölçekli bir mühendislik çizimidir.

\gorsel{img/tr-02-results.png}{Örnekte regresyon hızı yanma boyunca
azalırken (port genişledikçe kütle akısı düşer) port çapı büyür.}

\subsection{5. Adım — Dışa aktarın}

Sonuç panelindeki dışa aktarma düğmeleriyle şunları üretebilirsiniz:
teknik resim ve rapor PDF'leri, CAD dosyaları (STEP kesit eşleme dahil) ve
CSV/Excel veri tabloları. Çıktılar \texttt{Belgeler/HRMA} klasörüne yazılır.

% ================= 5. DİĞER MOTOR TİPLERİ =================
\section{Katı ve Sıvı Motorlar}

Katı ve sıvı sayfaları aynı deseni izler; yalnız girdiler tipe özeldir.

\begin{itemize}[leftmargin=1.4em]
  \item \textbf{Katı motor (\texttt{/solid}):} Grain geometrisi (BATES,
  finocyl, slotted vb.), segment sayısı, propelan ailesi (KNDX, KNSB, APCP
  ve dahası), balistik ve Monte Carlo belirsizlik çözümlemesi. Yanma hızı
  için \texttt{a} ve \texttt{n} katsayıları veri tabanından gelir.
  \item \textbf{Sıvı motor (\texttt{/liquid}):} Propelan çifti, besleme
  sistemi tipi ve \textbf{güç çevrimi} (basınç beslemeli, gaz jeneratörlü,
  tap-off, staged combustion, FFSC, expander). Çevrim güç dengesi çözücüsü
  pompa ve türbin güçlerinin kapanıp kapanmadığını kontrol eder. Ayrıca
  süperkritik CH\textsubscript{4}/H\textsubscript{2} rejeneratif soğutma ve
  gaz-gaz enjektör yolu vardır. Yanma termokimyası, motorun \emph{gerçek}
  oda basıncında, O/F oranında ve genişleme oranında RocketCEA ile çözülür.
\end{itemize}

% ================= 6. KAPSAM =================
\section{Kapsam ve Sınırlar (Mutlaka Okuyun)}

HRMA güçlü bir ön tasarım ve eğitim aracıdır, ancak sınırlarını bilmek
şarttır:

\begin{itemize}[leftmargin=1.4em]
  \item \textbf{Ön tasarım seviyesi.} Modeller kapalı-biçim ve 1B
  bağıntılardır. Turbopompa haritaları, kavitasyon, rotordinamik, yatak
  ömrü, başlatma/durdurma geçici rejimleri ve 2B/3B eşlenik ısı transferi
  \emph{yoktur}.
  \item \textbf{Doğrulama (validation) sınırlı.} Sıvı vakum Isp'si, gerçek
  çalışma noktası ve yayımlanmış geometri verildiğinde 14 motorluk bir
  örneklemde $\sim\%1.2$ medyan hatayla yeniden üretiliyor. Bu, ``bilinen
  noktada Isp'yi iyi tahmin eder'' demeye yeter; ``sıfırdan bilinmeyen bir
  motorun uçuş performansını bu hatayla tahmin eder'' demeye \emph{yetmez}.
  Hibrit tarafında regresyon hızı ve oda basıncı hataları çok daha büyüktür
  (\%13--35). Güncel sayılar için \texttt{docs/VALIDATION\_STATUS.md}.
  \item \textbf{Uçuşa elverişlilik değil.} HRMA, uçacak bir motorun nihai
  tasarımını, güvenlik kararını veya kalifikasyonunu tek başına
  desteklemez. Bunlar için yorulma/creep analizi, fiziksel test ve resmi
  bir V\&V süreci gerekir.
\end{itemize}

\begin{uyarikutu}
\small\color{hrmadark}
\textbf{Altın kural.} HRMA bir başlangıç noktasıdır. Kritik her sonucu
bağımsız bir kodla (CEA / RPA / openMotor) ve mümkünse statik ateş testiyle
doğrulayın. Programı kullanmayı bilmek, tek başına roket motoru mühendisi
olmak değildir.
\end{uyarikutu}

% ================= 7. GÜNCELLEME =================
\section{Otomatik Güncelleme}

HRMA açılışta GitHub Releases üzerinden yeni sürüm olup olmadığını denetler.
Yeni sürüm varsa bir pencere çıkar: \textbf{Şimdi güncelle}, \textbf{Daha
sonra} veya \textbf{Bu sürümü atla}. ``Şimdi güncelle'' derseniz uygulama
kurulumu indirir, kendini kapatır, sessizce kurar ve otomatik yeniden açılır.

\begin{notkutu}
\small\color{hrmadark}
\textbf{Sürüm 2.6.1 notu.} macOS'ta güncelleme kurulumu birkaç dakika
sürebilir ve bu sırada bir bildirim gösterilir; kurulum bitene kadar
uygulamayı elle açmayın. Ağınız GitHub'ın istek sınırına takılırsa
(paylaşımlı ağlarda olabilir) program artık normal sürüm sayfalarına
düşerek güncellemeyi yine de bulur.
\end{notkutu}

\vspace{1cm}
\begin{center}
\color{hrmagray}\small
HRMA v2.6.1 — Roket Motoru Analiz Paketi\\
\href{https://github.com/berketez/HRMA}{github.com/berketez/HRMA}
\end{center}

\end{document}
